% methods_materials.tex
\subsection{Materials}

The demonstration uses the 50-item International Personality Item Pool (IPIP) Big-Five Factor Markers, a public-domain inventory of ten short, first-person items (e.g., ``I am the life of the party.'') per domain: Extraversion, Emotional Stability, Agreeableness, Conscientiousness, and Openness/Intellect. The set operationalizes the five-factor structure \citet{goldberg-1990} demonstrated across near-comprehensive pools of English trait adjectives---an ideal test bed for a semantic method, pairing one of psychometrics' most replicated target structures with brief, behaviorally concrete items whose meaning a language model can plausibly capture. The inventory ships in the package as \texttt{data(big5)}: item texts, codes (E1--O50), precomputed item embeddings (see the Embeddings subsection), the official $\pm 1$ scoring key---under which \valNRev{} of the 50 items are reverse-keyed, with Emotional Stability keyed in the Neuroticism direction---and the theoretical domain labels used in all analyses, tables, and figures below, which name the second domain Neuroticism (its keyed direction) and the fifth Openness (its short form). The scoring metadata enter only the explicitly keyed procedures (the \texttt{atomic\_reversed} encoding and the scoring of the human responses); the primary similarity encoding is keying-free.

The human benchmark is the Open-Source Psychometrics Project's ``Big Five Personality Test'' raw-data release \citep{openpsychometrics-2018}, an open administration of the same 50 items on a 5-point Likert scale in which 0 codes a skipped item. Of its \valNHumanRaw{} response records we removed every record with a missing response or one outside the 1--5 range, retaining $N = \valNHuman$ respondents; the same release has served as a validation corpus in earlier embedding work \citep{wulff-mata-2025,jung-seo-2025}. These responses serve \emph{only} as the validation target---a conventional exploratory factor analysis (EFA) of their correlations against which the semantic solution is compared---and never enter the semantic analyses.
